\documentclass[11pt]{article}
\usepackage[T1]{fontenc}
\usepackage{lmodern,amsmath,amssymb,amsthm}
\usepackage[margin=1in]{geometry}
\usepackage{microtype}
\usepackage[colorlinks=true,linkcolor=blue,citecolor=blue,urlcolor=blue]{hyperref}
\newcommand{\Q}{\mathbb Q}
\newcommand{\R}{\mathbb R}
\newcommand{\N}{\mathbb N}
\newcommand{\SOS}{\sum F[X]^2}
\newtheorem{theorem}{Theorem}
\newtheorem{lemma}{Lemma}
\title{A fixed moment separator without a rational\\ normalized sum-of-squares certificate}
\author{Alec Kriebel\\
  \small Independent researcher\\
  \small \href{https://orcid.org/0009-0001-9320-500X}{ORCID: 0009-0001-9320-500X}}
\date{October 6, 2026}
\begin{document}
\maketitle
\begin{abstract}
We apply Scheiderer's classical rational quartic, which is a sum of real
squares but not rational squares, to the fixed-polynomial certificate question
in Oberwolfach Report 14/2023. On the rational unit ball in three variables,
an explicit rational degree-four moment vector has mass one and annihilates
a rational polynomial \(p\) with minimum one. The identity
\(p\in1+Q_{\R}(g)\) has a real certificate, whereas
\(p\notin1+Q_{\Q}(g)\), even allowing arbitrary finite multiplier degrees.
Here rational certificates mean rational polynomial square factors.
The all-degree obstruction is a classical local leading-form argument.
The moment matrix is not positive semidefinite, another polynomial gives a
rational certificate for the same vector, and every rational rescaling
\(cp\), \(c>1\), has a rational normalized certificate.
Thus the example exhibits the phenomenon for one fixed boundary separator;
it does not address the companion paper's strict-margin regime.
\end{abstract}

\section{The question and certificate convention}
The end of Naldi's contribution, joint work with Henrion and Safey El Din,
in Oberwolfach Report 14/2023 \cite[pp.~802--804]{OWR} asks whether a
rational polynomial certifying truncated-moment nonrepresentability can have
no rational certificate for its fixed membership in \(1+Q(g)\).
We exhibit this phenomenon under the rational-square convention below.
The quartic, its irrationality mechanism, and the local obstruction used
here are established results. The purpose of this note is their explicit
application to the printed moment question and clarification of its
normalization boundary.

For \(F=\Q\) or \(\R\), \(X=(x,y,z)\), and \(g\in F[X]\), define
\[
 Q_F(g)=
 \left\{\sum_{i=1}^r a_i^2+g\sum_{j=1}^s b_j^2:
 r,s\geq0,\quad a_i,b_j\in F[X]\right\}.
\]
All sums are finite; there is no degree bound on the factors.
A \emph{rational certificate} means that these factors lie in \(\Q[X]\).
This is equivalent to rational positive-semidefinite Gram data
\cite[Lemma~1.6]{Gram}. It is stronger than requiring rational coefficients
only in the real-SOS multiplier polynomials. The printed OWR question does
not formally specify this distinction; its Scheiderer reference provides
the rational-SOS context. We state our convention explicitly rather than
attribute an unstated formal definition to that source.

For \(m=(m_\alpha)_{|\alpha|\leq4}\), let
\[
 L_m\left(\sum_{|\alpha|\leq4}p_\alpha X^\alpha\right)
       =\sum_{|\alpha|\leq4}p_\alpha m_\alpha .
\]
A representing measure is a nonnegative Borel measure supported on
\(K=\{X\in\R^3:g(X)\geq0\}\), with the specified moments.

\section{Explicit example and proof}
Use the classical Scheiderer quartic
\cite[Theorem~2.1 and Example~2.8]{Sch16}:
\begin{equation}\label{eq:f}
\begin{split}
 f={}&x^4+xy^3+y^4-3x^2yz-4xy^2z\\
     &+2x^2z^2+xz^3+yz^3+z^4 .
\end{split}
\end{equation}
Set
\begin{equation}\label{eq:data}
 g=1-x^2-y^2-z^2,\qquad p=1+f,\qquad
 m_{000}=1,\quad m_{400}=-1,\quad
 m_\alpha=0\ \text{otherwise}.
\end{equation}
There are \(\binom{7}{3}=35\) moment entries.

\begin{theorem}\label{thm:main}
For the data \eqref{eq:data}, \(K\) is the nonempty compact unit ball,
and \(Q_{\Q}(g)\) is Archimedean. The vector \(m\) has no representing
measure, even on \(\R^3\), and \(L_m(p)=0\). Moreover,
\[
 \min_Kp=1,\qquad
 p\in1+Q_{\R}(g),\qquad p\notin1+Q_{\Q}(g).
\]
The last nonmembership holds for every finite choice of polynomial
square factors, of arbitrary degrees.
\end{theorem}

We recall the credited quartic argument to make the obstruction checkable.
No new quartic or Galois construction is claimed.

\begin{lemma}[Scheiderer's quartic]\label{lem:quartic}
The polynomial \(f\) in \eqref{eq:f} is a sum of real polynomial squares
and is not a sum of rational polynomial squares.
\end{lemma}
\begin{proof}
Let \(h(t)=t^4-t+1\). It has no real roots: for \(t\leq0\) or
\(t\geq1\) its value is at least one, and for \(0<t<1\) both
\(t^4\) and \(1-t\) are positive. Modulo \(2\), \(h\) is irreducible:
it has no linear factor and is not divisible by the sole irreducible
monic quadratic \(t^2+t+1\). Modulo \(3\),
\[
 h(t)=(t+1)(t^3-t^2+t+1),
\]
where the cubic is irreducible and the factors are distinct.
The usual factorization-cycle theorem gives a \(4\)-cycle and a \(3\)-cycle
in the Galois group. Its order therefore is \(12\) or \(24\);
the index-two subgroup \(A_4\) contains no \(4\)-cycle. Thus the group is
\(S_4\).

For the four roots \(\alpha_i\), put
\(\ell_i=x+\alpha_i y+\alpha_i^2z\).
The norm identity \(f=\prod_i\ell_i\) follows by expansion, or by the
determinant of multiplication by \(x+ty+t^2z\) in \(\Q[t]/(h)\).
Every three \(\ell_i\) are independent by the Vandermonde determinant.
Complex conjugation pairs the roots. Choosing one root from each pair,
\(f=|\ell_1\ell_2|^2\), after relabeling, is a sum of two real
quadratic squares.

Suppose \(f=\sum_j q_j^2\) with rational homogeneous quadratics \(q_j\).
The intersection of each conjugate pair of projective lines
\(\{\ell_i=0\}\) is a real projective point, since its one-dimensional
complex vector space is conjugation-stable. At those points \(f=0\),
so every \(q_j\) vanishes. Rational coefficients and the \(S_4\) action
then force vanishing at all six pairwise line intersections.
On each of the four lines there are three distinct such points.
A binary quadratic vanishing at three distinct projective points
vanishes identically on that line. Hence all four distinct \(\ell_i\)
divide each \(q_j\), impossible for a nonzero quadratic.

An inhomogeneous rational-square representation gives no escape.
The highest homogeneous square terms cannot cancel over \(\R\),
so all factors have degree at most two. The constant part of \(f\)
forces all constant factors to vanish; its degree-two part then
forces their linear parts to vanish. The factors would be precisely
the homogeneous quadratics just excluded.
\end{proof}

For an explicit real certificate, let \(\beta\in(-2,-1)\) satisfy
\(\beta^3-4\beta-1=0\). Scheiderer's displayed identity is
\begin{equation}\label{eq:beta}
\begin{gathered}
 U=2x^2+\beta y^2-yz+(2+\beta^{-1})z^2,\\
 V=2xy-\beta^{-1}y^2+2\beta^{-1}xz+\beta yz-z^2,\\
 4f=U^2-\beta V^2.
\end{gathered}
\end{equation}
Existence of this negative root follows from opposite signs at \(-2,-1\).
Thus \(f=(U/2)^2+(\sqrt{-\beta}\,V/2)^2\).
The supplement checks \eqref{eq:beta} exactly after clearing denominators,
modulo the cubic.

\begin{proof}[Proof of Theorem~\ref{thm:main}]
The geometry of \(K\) is immediate, and its displayed rational ball
generator makes the module Archimedean.
The negative fourth moment \(m_{400}\) excludes every nonnegative
representing measure. The coefficient of \(x^4\) in \(f\) is one, so
\(L_m(p)=1-1=0\).
Lemma~\ref{lem:quartic} gives \(f\geq0\) and
\(p\in1+Q_{\R}(g)\). Since \(f(0)=0\), the minimum is one.

For the remaining assertion, suppose a finite rational identity
\[
 f=\sum_i a_i^2+(1-x^2-y^2-z^2)\sum_j b_j^2
\]
holds, with arbitrary factor degrees.
At the origin, both weights are one, so every \(a_i(0),b_j(0)\) vanishes.
The homogeneous degree-two part is then the sum of the squares of all
linear parts. It is zero, so these linear parts also vanish.
Writing \(a_i^{[2]},b_j^{[2]}\) for the homogeneous quadratic parts,
the degree-four part of the identity is exactly
\[
 f=\sum_i (a_i^{[2]})^2+\sum_j(b_j^{[2]})^2.
\]
The product of \(g-1\) with any \(b_j^2\) starts at degree six;
higher factor terms do not affect degree four. This contradicts
Lemma~\ref{lem:quartic}. No global degree bound or assertion that
higher weighted terms cannot cancel was used.
\end{proof}

This local argument is classical. Scheiderer's leading-form lemma and
positive-generator completion argument already give its mechanism
\cite[Lemma~1.1 and proof of Proposition~6.1]{Sch00}.
Benoist explicitly observes that this same quartic is not a sum of
squares in \(\Q[[x,y,z]]\) \cite[Remark~2.7]{Benoist}.
Alternatively, a rational ball-module identity would become a rational
formal-square identity after multiplying the \(b_j\) by the binomial
series \(\sqrt{1-x^2-y^2-z^2}\).
Nie previously applied the corresponding leading-form obstruction over
the reals to the homogeneous Motzkin form on a ball
\cite[Example~5.3]{Nie}.

\section{Normalization, scope, and related work}
The following controls are part of the result, not optional qualifications.
First, the same \(m\) has the rationally certifiable separator
\[
 q=1+x^4,\qquad L_m(q)=0,\qquad q-1=(x^2)^2 .
\]
Second, Powers' rational Putinar theorem for a rational ball generator
\cite[Theorem~7]{Powers} gives \(p\in Q_{\Q}(g)\):
\(p\) itself is strictly positive on \(K\).
For every rational \(c>1\),
\[
 cp-1=(c-1)+cf\geq c-1>0 \quad\hbox{on }K,
\]
so the same theorem gives \(cp\in1+Q_{\Q}(g)\), with
\(L_m(cp)=0\). The obstruction concerns \(c=1\). For \(c<1\),
\(cp(0)-1<0\), so even real normalized membership fails.
No degree bound for Powers' certificates is asserted.

Third, \(M_2(m)\), indexed by monomials of degree at most two, is not
positive semidefinite: its diagonal entry indexed by \(x^2\) is
\(m_{400}=-1\). The construction does not address a PSD-input
restriction. Finally, a merely rational-coefficient list of real-SOS
multipliers already exists, namely \(\sigma_0=f,\sigma_1=0\).
It is the rational square factors that are impossible.

The companion preprint of Henrion, Naldi, and Safey El Din
\cite[Definition~3, Corollary~2, and Remark~2]{TMP} distinguishes a
strict-margin sufficient regime, \(\min_Kp>1\). Our minimum-one example
addresses the literal fixed-polynomial OWR question and is outside that
corollary's regime. It supplies neither an algorithmic claim nor a
minimal-relaxation-degree claim, and does not settle general
Archimedean descent variants.

The same quartic's irrational Gram points were studied by Henrion,
Naldi, and Safey El Din \cite[\S6.2]{LMI}. Naldi and Sinn already used it
to construct a strongly infeasible rational semidefinite program without
a rational PSD separating certificate \cite[Example~3.8]{NS}.
That matrix-certificate application is substantial prior work; the
present note uses the fixed normalized moment identity and arbitrary
polynomial multiplier degrees instead. Our contribution is this
explicit application and its scope clarification, rather than a new
irrationality mechanism.

A dated bounded priority note is included in the supplement. It records
the primary passages inspected and exact access limits, including the
unexamined remainder of the 2026 corrected journal proof on descent
\cite{Descent}, whose full body was unavailable. Its inspected
univariate, strict-positive, and zero-dimensional-support scopes do not
conflict with this example. The module here is already rationally
Archimedean, and its support ideal is zero because \(K\) has interior.
The quotient \(\Q[X]/(Q_{\Q}(g)\cap -Q_{\Q}(g))\) therefore has
dimension three. No statement about the unavailable Appendix~B
theorem is inferred. We make no worldwide-firstness or current-openness
assertion.

\paragraph{Verification and disclosure.}
The portable supplement contains this manuscript source, exact rational
data, a standard-library checker and its full expected output, scope and
priority notes, and a hash manifest. Its 2,059 finite checks cover the
norm identity, finite-field factorizations, real-SOS identity, moment
data, and representative leading-form/incidence controls. The proofs
above, not finite searches, establish the Galois and all-degree
nonmembership assertions. AI tools were used extensively in developing
the application, literature comparison, exact computations, writing,
and adversarial checking. This is an unrefereed preprint; those checks
do not constitute independent human peer review or formal
proof-assistant certification.

\begin{thebibliography}{10}
\bibitem{OWR}
S. Naldi, On Algebraic Certificates for the Truncated Moment Problem
(joint work with D. Henrion and M. Safey El Din),
in \emph{Real Algebraic Geometry with a View toward Koopman Operator Methods},
Oberwolfach Report 14/2023, pp.~802--804.
\href{https://doi.org/10.4171/OWR/2023/14}{doi:10.4171/OWR/2023/14}.

\bibitem{TMP}
D. Henrion, S. Naldi, and M. Safey El Din,
\emph{Algebraic certificates for the truncated moment problem},
arXiv:2302.06927v1 (2023).
\url{https://arxiv.org/abs/2302.06927v1}.

\bibitem{Sch16}
C. Scheiderer, Sums of squares of polynomials with rational coefficients,
\emph{J. Eur. Math. Soc.} 18 (2016), 1495--1513.
\href{https://doi.org/10.4171/JEMS/620}{doi:10.4171/JEMS/620}.

\bibitem{Sch00}
C. Scheiderer, Sums of squares of regular functions on real algebraic varieties,
\emph{Trans. Amer. Math. Soc.} 352 (2000), 1039--1069.
\href{https://doi.org/10.1090/S0002-9947-99-02522-2}{doi:10.1090/S0002-9947-99-02522-2}.

\bibitem{Benoist}
O. Benoist, On the bad points of positive semidefinite polynomials,
\emph{Math. Z.} 300 (2022), 3383--3403.
\href{https://doi.org/10.1007/s00209-021-02804-9}{doi:10.1007/s00209-021-02804-9}.

\bibitem{Nie}
J. Nie, \emph{An Exact Jacobian SDP Relaxation for Polynomial Optimization},
author manuscript dated May 26, 2011, Example~5.3, p.~21.
\url{https://mathweb.ucsd.edu/~njw/PUBLICPAPERS/JacSdp.pdf}.

\bibitem{LMI}
D. Henrion, S. Naldi, and M. Safey El Din,
Exact algorithms for linear matrix inequalities,
\emph{SIAM J. Optim.} 26 (2016), 2512--2539.
\href{https://doi.org/10.1137/15M1036543}{doi:10.1137/15M1036543}.

\bibitem{NS}
S. Naldi and R. Sinn, Conic programming: Infeasibility certificates and
projective geometry, \emph{J. Pure Appl. Algebra} 225 (2021), 106605.
\href{https://doi.org/10.1016/j.jpaa.2020.106605}{doi:10.1016/j.jpaa.2020.106605}.

\bibitem{Powers}
V. Powers, Rational certificates of positivity on compact semialgebraic sets,
\emph{Pacific J. Math.} 251 (2011), 385--391.
\href{https://doi.org/10.2140/pjm.2011.251.385}{doi:10.2140/pjm.2011.251.385}.

\bibitem{Gram}
L. Chua, D. Plaumann, R. Sinn, and C. Vinzant,
\emph{Gram Spectrahedra}, arXiv:1608.00234v1 (2016), Lemma~1.6.
\url{https://arxiv.org/abs/1608.00234v1}.

\bibitem{Descent}
M. K. Keshari, D. Ojha, and N. S. Patra,
\emph{Descent problem for certificate of non-negativity on semi-algebraic sets},
arXiv:2508.07060v1 (2025); corrected journal proof (2026),
\href{https://doi.org/10.1016/j.indag.2026.09.003}{doi:10.1016/j.indag.2026.09.003}.
The journal full body and Appendix~B statement were not audited.
\end{thebibliography}
\end{document}
